%------------------------------------------------------------------------------
%------------------------------------------------------------------------------
%
%   Script Name : preamble.tex
%        Author : Ecometer s.n.c.
%          Date : 2024-09-30
%   Description : LaTeX preamble file
%
%------------------------------------------------------------------------------
%------------------------------------------------------------------------------

%------------------------------------------------------------------------------
% DOCUMENT CLASS
%------------------------------------------------------------------------------
\documentclass[a4paper,10pt]{scrartcl} % article (KOMA-Script)

%------------------------------------------------------------------------------
% FONT
%------------------------------------------------------------------------------
%% font 'Roboto Light'
\usepackage[sfdefault,light]{roboto}
\usepackage[T1]{fontenc}

%------------------------------------------------------------------------------
% PACKAGES
%------------------------------------------------------------------------------
\usepackage[utf8]{inputenc} % encoding used in input file if it is not ascii
\usepackage[italian]{babel} % convenzioni tipografiche italiane
\usepackage[left=1cm, right=1cm, top=0.5cm, bottom=3.5cm, includeheadfoot]{geometry} % permette la modifica della gabbia del documento
\usepackage[pdftex]{graphicx} % inclusione di oggetti grafici esterni (\includegraphics)
\usepackage{color} % per testo colorato

\usepackage{geometry} % per gestire i margini delle minipage
\geometry{
    a4paper,
    total={170mm,257mm},
    left=5mm,
    right=5mm,
}

\usepackage{enumitem} % modificare gli ambienti enumerate, itemize e description
\usepackage{fancyhdr} % facilities per la costruzione di headers e footers
\usepackage{lastpage} % allow to refer to the last page of a document: \pageref{LastPage}
\usepackage{subfigure} % support for the inclusion of small ???gures and tables
\usepackage{amssymb} % symbols
\usepackage{typearea} % setting typearea for KOMA-script
\usepackage{verbatim} % contains a reimplementation of the verbatim environment
\usepackage{booktabs} % per tabelle
\usepackage{array} % per tabelle
\usepackage{longtable} % per tabelle su più pagine
\usepackage{tabu} % Flexible LATEX tabulars
\usepackage{wasysym} % per disegnare pallini vuoti/pieni \Circle \CIRCLE
\usepackage{tikz} % create graphic elements
\usepackage{tabularx} % Tabulars with adjustable-width columns
\usepackage{tasks} % Horizontally columned lists
\usepackage{pifont} % Access to PostScript standard Symbol and Dingbats fonts
\usepackage{xcolor} % per colorare le linee dell'header e del footer
\usepackage{adjustbox} % Graphics package-alike macros for “general” boxes
\usepackage{hhline} % Better horizontal lines in tabulars and arrays
\usepackage{colortbl} % colori linee tabelle
\usepackage{textcomp} % per � simbolo

% per aggiungere hyperlink
\usepackage{hyperref}
\hypersetup{
    colorlinks=true,
    linkcolor=blueCF,
    filecolor=blueCF,
    urlcolor=blueCF,
}

\usepackage{multirow} % Create tabular cells spanning multiple rows
\usepackage{datetime2} % per inserire data e ora

\usepackage{multicol} % Intermix single and multiple columns
\setlength{\columnsep}{0.5cm}

\usepackage[font=small,labelfont=bf]{caption} % Required for specifying captions to tables and figures
\usepackage{ragged2e} % Alternative versions of “ragged”-type commands
\usepackage{soul} % Hyphenation for letterspacing, underlining, and more
\usepackage{float} % Improved interface for floating objects

%------------------------------------------------------------------------------
% RENEWCOMMANDS BY USER
%------------------------------------------------------------------------------
\renewcommand{\arraystretch}{1.23} % this reduces the vertical spacing between rows
